\section{讨论、有效性契约与审查门禁协议}
\label{sec:discussion}

\subsection{物理约束预测模型的四项实证边界条件}
系统审计、代码验证与前瞻性消融实验共同指向四个边界条件，决定物理先验何时有效、何时中性、何时有害。

\paragraph{样本受限工况下的防御性收益与先验误设脆弱性}
在数据受限工况（Penmanshiel，$N_{\text{train}}=1,583$），$\lambda=1.0$ 下校准的功率曲线先验将归一化非对称调度成本从 $0.09881 \pm 0.00027$ 稳步降至 $0.09760 \pm 0.00044$（降低 $-1.22\%$），提供稳健的下行风险保护。相比之下，尺度对齐的打乱先验诱发成本恶化至 $0.16611 \pm 0.00335$（激增 $+68.1\%$）。两项实证结果的不对称性揭示了强物理正则化的非对称风险—收益包线：强归纳护栏在抑制自由假设空间的同时，成倍放大了先验形式化中的潜在误设误差。跨场站对比无法剥离这一效应，因为场站间的地形特征与机型差异会与样本量同时变化。数据充裕时（Kelmarsh，$N_{\text{train}}=5,098$），无约束 DLinear（$0.08812$）、真实物理（$0.08791$）与打乱先验（$0.08801$）在数值上高度密集，表明模型容量已饱和假设空间。

\paragraph{低惩罚尺度表现为非特异性数值正则化}
现有文献常在损失函数中配置较小的物理惩罚系数（$\lambda \in [10^{-4}, 10^{-2}]$）。受控消融实验证实：在 $\lambda = 0.01$ 下，真实物理先验在 Kelmarsh 取得的调度成本（$0.08810 \pm 0.00009$）与打乱物理关系的伪先验（$0.08801 \pm 0.00008$）高度重合且种子方差重叠。在弱权重区间下，损失惩罚项在实证上表现为非特异性的数值平滑，尚未显现空气动力学因果引导特征。确立严谨的物理特异性，需引入形式与权重匹配的伪打乱阴性对照并在 $\lambda \ge 0.1$ 区间实施敏感性审计。

\paragraph{训练损失物理代理时效错配诱发的成本单调膨胀}
当物理约束形式化依赖外部气象代理计算损失惩罚项（如通过风速查询经验功率曲线）时，在缺乏前向数值天气预报的离线训练场景中，研究者常倾向于使用滞后历史气象观测代偿。对滞后步数 $k \in \{6, 12, 18, 36\}$（对应 1 至 6 小时时效错配）的参数扫描揭示，三座风场的调度性能随着滞后增加呈现严格单调恶化。在 Hill of Towie，决策风险从 $k=6$ 时的 $0.09022$ 持续上升至 $k=36$ 时的 $0.09257$（恶化 $+2.6\%$），在约 $k=12$ 时即劣于无约束基准（$0.08924$）。我们特别强调，该实验刻画的是离线训练目标中气象代理时效错配的敏感性，而非测试推理阶段的遥测延迟（因为骨干网络推理时仅读取历史功率序列）。然而，它表明若离线训练目标依赖滞后气象代理构造正则项，时效错配将迅速削弱先验的归纳价值，构成依赖代理变量的约束架构的实证应用局限。

\paragraph{评价指标脱钩、格式化远期与不平衡核算及退化零申报诊断探针}
在格式化批发电力核算模型中，发电机组总收益由远期合约持仓与实时平衡偏差结算共同决定：$R_{\text{total}} = P_{\text{F}} \cdot q + P_{\text{imb}} \cdot (y - q)$，其中 $q$ 为核算模型中的远期签约或申报电量头寸。实证评估若剥离远期履约收益并孤立计算未对冲的平衡偏差现金流，所度量的实为不平衡租金而非预测技能。作为一项不平衡结算与决策风险的分析性诊断探针，退化零申报基准（$q \equiv 0$）构成了严格的数学归谬（\textit{reductio ad absurdum}）：它完全不产生任何预测，仅通过将全部物理发电量转化为未履约余量并按 $P_{\text{imb}}$ 结算，即可在三座风场凭空捕获 22.3 万至 24.3 万英镑/台$\cdot$年的未对冲正向现金流。这从数理上证实了孤立未对冲的平衡现金流与预测技能完全脱钩；评估真实的商业价值必须联合建模远期合约收益（$P_{\text{F}} \cdot q$）与实时不平衡结算。实际运维评估必须采用与市场非对称惩罚比率相匹配的 Pinball 损失或耦合远期与平衡结算的综合核算。

\subsection{外部有效性边界契约}
\label{sec:boundary_contract}
为确立清晰的工程适用范围，本节将实证发现归纳为四项运行有效性边界条件：

\paragraph{边界一：样本规模与假设空间饱和}
场站内学习曲线证实，在我们测试的配置下，物理约束的有效收益主要集中于数据受限工况（$N_{\text{train}} \le 1,000$，图~\ref{fig:within_site_learning_curve}）。当模型训练窗口扩增至 $N_{\text{train}} = 5{,}000$ 时，假设空间饱和将真实先验与无约束基准的差距缩小至 0.33\% 以内（Kelmarsh），且无约束模型略微反超，而打乱先验在所有样本规模下均造成严重破坏。连续退化测试进一步确立了狭窄的工程容错窗口（Penmanshiel 上临界阈值 $\alpha^* \approx 0.07$，图~\ref{fig:corruption_gradient}）。机组气动老化折损超出 10\% 或测风仪标定偏差超出 $\pm 5\%$ 时，标称物理先验即翻转为运维负债。

\paragraph{边界二：代理时效对齐与观测陈旧度}
对于依赖代理变量的物理约束架构，先验有效性严格受限于气象代理的时间对齐（$k \le 6$ 步）。在代理时效错配跨度超过 1 小时的工况下，三座风场均观察到成本的严格单调上升（Hill of Towie 在 $k=36$ 时成本上涨 $+2.6\%$），性能显著劣于无约束基准。在训练目标中以滞后历史代理替代同步气象数据的做法会诱发显著的性能退化。

\paragraph{边界三：市场结构与格式化不平衡结算核算}
本研究所揭示的经济实证结论，本质上作为不平衡结算与决策风险的诊断探针，严格受限于具备非对称短缺惩罚比率（$r \ge 1.0$）与半小时时段不平衡暴露聚合的英国 BSC 单一现金流结算机制。单边不平衡现金流并不代表完整的商业价值，且不外推至对称单一结算体系或固定上网电价（FIT）机制。评价者不应将孤立未对冲的平衡现金流指标外推为衡量预测模型的运维实用准则。

\paragraph{边界四：软件保真度与计算图审计范畴}
软件层面的形式化审计结论严格适用于 PCT 开源计算图（commit \texttt{65941ec}）及受评的线性（DLinear）与块状注意机制（PatchTST）架构。虽然方向颠倒的损失散度在 PCT 的 AST 节点上得到严格证实，但对于未开源代码的专有模型，必须经过独立的静态计算图审计后方可确认其物理保真度。

\subsection{两项预先设立假设的证伪结果}
本研究记录了在分析前预先设立并在实证中被严格证伪的两项假设。张量代数推导证伪了关于 PCT 乘积核 KDE 导致边缘独立的初步猜想，严格证明其二维耦合在代数展开中得以保留。对称地，电力市场诊断探针证伪了“校准后反转源自物理正则化”的假设，证明现金流逆转完全受控于非对称结算结构。如实记录假设证伪过程，使整个审计追踪具备可复现的认识论严谨性。

\subsection{REFORMS-PCWPF 方法学报告清单规范}
\label{sec:reforms_pcwpf}
为将实证审计发现转化为长效的方法学报告规范，表~\ref{tab:reforms_checklist} 提出了面向未来物理约束风电预测研究的五道推荐验证门禁：

\begin{table*}[!t]
\centering
\footnotesize
\renewcommand{\arraystretch}{0.88}
\setlength{\tabcolsep}{3.5pt}
\caption{面向可信物理约束新能源预测的 REFORMS-PCWPF 报告规范清单（五大模块化合规门禁）。}
\label{tab:reforms_checklist}
\begin{threeparttable}
\begin{tabularx}{\textwidth}{p{1.2cm}p{2.6cm}>{\raggedright\arraybackslash}X>{\raggedright\arraybackslash}Xp{3.2cm}}
\toprule
\textbf{模块编号} & \textbf{治理核心维度} & \textbf{强制性审计实施协议} & \textbf{常见方法学缺陷 / 反模式} & \textbf{必须提交的验证工件} \\
\midrule
\textbf{M1} & \textbf{先验隔离消融} & 训练完全相同的网络骨干，仅在相同的随机种子下对比 $\lambda=0$ 与 $\lambda>0$。 & 混淆模块消融（直接砍掉整个 GNN 或注意力子块）与物理领域知识。 & 参数量对齐且具备多随机种子方差的 $\lambda=0$ vs. $\lambda>0$ 对比表。 \\
\midrule
\textbf{M2} & \textbf{领域特异性对照} & 评估保持损失形式、模型架构与正则化权重 $\lambda$ 完全相同、但打乱物理映射的伪先验。 & 将低惩罚尺度（$\lambda \le 0.01$）下的非特异性数值正则化误判为物理推理。 & 伪打乱对照性能曲线与实证特异性比率（$p$ 值假设检验显著性）。 \\
\midrule
\textbf{M3} & \textbf{代理陈旧度压力测试} & 若约束依赖外部物理输入（如风速），扫描静态代理陈旧度 $\tau \in \{1\text{h}, 2\text{h}, 3\text{h}, 6\text{h}\}$ 评估在缺失前瞻数值天气预报时的静态代理陈旧度影响。 & 默认在测试推理期能够无延迟、无噪声地获取理想实时输入。 & 展现有效性受损临界时域的代理陈旧度成本膨胀曲线。 \\
\midrule
\textbf{M4} & \textbf{指标解耦与结算核算} & 综合汇报对称点误差 (RMSE)、非负非对称决策风险 (Pinball)，测试退化零申报反例。 & 孤立报告未对冲的单边偏差现金流，误将系统偶然的价格尖峰低估当做技能。 & 耦合远期合同持仓与平衡现金流的解耦市场评估。 \\
\midrule
\textbf{M5} & \textbf{计算图保真度} & 审查软件中的损失函数实现是否在参数传入方向和归一化上与发表数学公式一致。 & 非对称散度（如 PyTorch \texttt{F.kl\_div}）参数顺序颠倒导致反向优化。 & 静态代码审查报告以及验证损失张量梯度计算方向的单元测试脚本。 \\
\bottomrule
\end{tabularx}
\begin{tablenotes}[flushleft]
\footnotesize
\item \textbf{注记}：基于对 46 篇范围界定文献的系统审计推导确立，严格对齐跨学科 REFORMS 规范 \cite{kapoor2024reforms} 与 DOME 指南 \cite{walsh2021dome}。我们建议将模块 M1--M5 作为未来物理约束新能源预测研究的推荐报告规范清单。
\end{tablenotes}
\end{threeparttable}
\end{table*}


\begin{itemize}
\item \textbf{门禁 M1（参数隔离消融）：} 建议评审避免接受将剔除大型神经网络子模块（如多头注意力、图卷积）的实验标称为“物理消融”，强制要求在参数对齐的网络容量下开展对照（$\lambda=0$ vs. $\lambda>0$），以剥离模型容量干扰。
\item \textbf{门禁 M2（伪打乱阴性对照）：} 可信物理有效性论证应包含形式与权重匹配的置换伪打乱对照，以严格区分空气动力学先验知识与非特异性数值平滑效应。
\item \textbf{门禁 M3（运行代理陈旧度压力测试）：} 约束评估应避免假定零延迟未来气象输入；在缺少前瞻 NWP 预报时，应主动汇报在真实代理陈旧度（$k \in \{6, \dots, 36\}$ 步）下的性能响应轨迹。
\item \textbf{门禁 M4（远期与平衡结算耦合评估）：} 行业评估应避免依赖孤立未对冲的平衡单边现金流，推荐采用与市场非对称惩罚比率匹配的 Pinball 损失或耦合远期持仓与平衡现金流的综合核算（$P_{\text{F}} \cdot q + P_{\text{imb}} \cdot (y - q)$）。
\item \textbf{门禁 M5（静态计算图与梯度核验）：} 鼓励作者开源可执行的训练脚本与计算图定义，建议评审核验损失函数的张量参数传递顺序与约简语义是否与论文数学公式严格相符。
\end{itemize}
